% SOURCES: figures/*.png (Mermaid architecture diagrams); docs/0.2.0v/architecture.md;
%          docs/0.2.0v/element-model.md; src/document/types.ts (ground truth for the
%          class diagram); src/export/index.ts (9-stage pipeline); src/draw/* (flow).
% The three formal UML diagrams (class / use-case / sequence) are drawn natively in
% TikZ (styles defined in main.tex), so no external diagram tool is required.

\chapter{System Analysis}
\label{ch:analysis}

This chapter analyses the system: first the layered architecture and data flow,
then the three formal UML views required for the design --- a \emph{class diagram}
of the Document Model, a \emph{use-case diagram} of the actors and their goals, and
\emph{sequence diagrams} of the two most characteristic interactions
(draw-to-create and export).

\section{Architectural Overview}
Draw-to-Web is organised into layers with a strictly one-directional data flow: the
user interface writes to the store through operations, the Document Model is the
single source of truth, and the output layer reads the document to generate and
export. Nothing flows the other way --- the canvas is a \emph{rendering} of the
document, never the model itself. Figure~\ref{fig:architecture} shows the layered
architecture across the Electron main, preload, and renderer processes, and
Figure~\ref{fig:dataflow} shows the one-directional data flow that this layering
enforces.

\begin{figure}[H]
  \centering
  \includegraphics[width=0.92\linewidth]{01-architecture.png}
  \caption{Layered architecture: renderer (UI $\rightarrow$ store $\rightarrow$ generator) over a sandboxed preload bridge and an Electron main process.}
  \label{fig:architecture}
\end{figure}

\begin{figure}[H]
  \centering
  \includegraphics[width=0.92\linewidth]{02-dataflow.png}
  \caption{One-directional data flow. The UI dispatches operations to the store; the generator and export pipeline read the resulting document.}
  \label{fig:dataflow}
\end{figure}

\begin{figure}[H]
  \centering
  \includegraphics[width=0.86\linewidth]{04-component-hierarchy.png}
  \caption{Renderer component hierarchy.}
  \label{fig:components}
\end{figure}

\begin{figure}[H]
  \centering
  \includegraphics[width=0.78\linewidth]{05-state-stores.png}
  \caption{State stores (document, history, session) and their relationships.}
  \label{fig:stores}
\end{figure}

\section{Class Diagram (Document Model)}
The Document Model is the heart of the system and therefore its class model. A
single \texttt{Document} aggregates the design tokens, the SEO configuration, the
runtime flags, the author settings, the variables, and the asset manifest, plus one
element \texttt{tree}. Every renderable element is an \texttt{ElementNode} --- an
abstract base (all elements share an identity, a per-breakpoint \texttt{style}, a
\texttt{states} map, and an optional animation) specialised into eight concrete
primitive types. Only \texttt{ContainerNode} owns children, and its children are
themselves \texttt{ElementNode}s, which is the recursive composition that makes the
document a tree. Complex shapes (hero, cards grid, navigation) are \emph{not} new
classes: they are presets that compose these same primitives.
Figure~\ref{fig:class} shows the model, generated from
\texttt{src/document/types.ts}.

\begin{figure}[p]
  \centering
  \resizebox{\textwidth}{!}{%
  \begin{tikzpicture}[node distance=10mm and 6mm]
    % ---- Root document ----
    \node[umlclass, text width=42mm] (doc) {
      \textbf{Document}
      \nodepart{two}
      +version: DocumentVersion\\
      +meta: DocumentMeta\\
      +tokens: Tokens\\
      +tree: ElementNode\\
      +seo: SEOConfig\\
      +runtime: RuntimeFlags\\
      +variables: DocumentVariables\\
      +settings: DocumentSettings\\
      +assets: AssetManifest
    };

    % ---- Document-level value objects (right of doc) ----
    \node[umlvalue, right=20mm of doc.north east, anchor=north west] (tokens) {
      \textbf{Tokens}\nodepart{two}+color[\,]\\+spacing[\,]\\+fontSize[\,]\\+radius[\,]\\+shadow[\,]
    };
    \node[umlvalue, below=3mm of tokens] (seo) {
      \textbf{SEOConfig}\nodepart{two}+title, description\\+openGraph, twitter\\+jsonLd\\+favicon, robots
    };
    \node[umlvalue, below=3mm of seo] (runtime) {
      \textbf{RuntimeFlags}\nodepart{two}+themeToggle\\+scrollSpy\\+mobileNav, reveals\\(8 opt-in flags)
    };
    \node[umlvalue, below=3mm of runtime] (settings) {
      \textbf{DocumentSettings}\nodepart{two}+contrastTarget\\+defaultTheme\\+decorativeBackdrop
    };

    % ---- Abstract element ----
    \node[umlabstract, below=24mm of doc, text width=52mm] (en) {
      \textit{\textbf{ElementNode}}\, \textit{(abstract)}
      \nodepart{two}
      +id: ElementId\\
      +name, semanticRole\\
      +style: ResponsiveProperties[StyleBlock]\\
      +states: StatesMap\\
      +animation: AnimationSpec\\
      +hiddenAt[\,]
    };

    % ---- Element-level value objects (LEFT of en; kept compact so the recursive
    %      "children" edge has a clear corridor below them) ----
    \node[umlvalue, left=18mm of en.north west, anchor=north east] (style) {
      \textbf{StyleBlock}\nodepart{two}+width, height, padding\\+background[\,], border, radius\\+typography, shadows[\,]
    };
    \node[umlvalue, below=3mm of style] (states) {
      \textbf{StatesMap}\nodepart{two}+hover, focus-visible, active
    };

    % ---- Concrete subclasses: two rows of four ----
    \node[umlnode, below=18mm of en, xshift=-45mm] (c1) {\textbf{ContainerNode}\nodepart{two}type='container'\\+children[\,]\\+layout: LayoutConfig};
    \node[umlnode, right=6mm of c1] (c2) {\textbf{TextNode}\nodepart{two}type='text'\\+tag: TextTag\\+content};
    \node[umlnode, right=6mm of c2] (c3) {\textbf{ImageNode}\nodepart{two}type='image'\\+alt, assetId\\+sizesHint};
    \node[umlnode, right=6mm of c3] (c4) {\textbf{ButtonNode}\nodepart{two}type='button'\\+content\\+ariaLabel};
    \node[umlnode, below=8mm of c1] (c5) {\textbf{LinkNode}\nodepart{two}type='link'\\+href, target\\+rel};
    \node[umlnode, below=8mm of c2] (c6) {\textbf{IconNode}\nodepart{two}type='icon'\\+name, inlineSvg\\+decorative};
    \node[umlnode, below=8mm of c3] (c7) {\textbf{ListNode}\nodepart{two}type='list'\\+ordered\\+items[\,]};
    \node[umlnode, below=8mm of c4] (c8) {\textbf{DividerNode}\nodepart{two}type='divider'\\+orientation};

    % ---- Compositions from Document ----
    \draw[compose] (doc.east) -- (tokens.west) node[midway, above, font=\tiny]{tokens};
    \draw[compose] (doc.east) -- (seo.west);
    \draw[compose] (doc.east) -- (runtime.west);
    \draw[compose] (doc.east) -- (settings.west);
    \draw[compose] (doc.south) -- (en.north) node[midway, right, font=\tiny]{tree 1};

    % ---- Compositions on the element (to the left) ----
    \draw[compose] (en.west) -- (style.east) node[midway, above, font=\tiny]{style};
    \draw[compose] (en.west) -- (states.east);

    % ---- Generalisation (nested shared bus) ----
    \coordinate (pbus) at ($(en.south)+(0,-6mm)$);
    \draw[inherit] (pbus) -- (en.south);
    \draw (c1.north |- pbus) -- (c4.north |- pbus);
    \foreach \c in {c1,c2,c3,c4}{ \draw (\c.north) -- (\c.north |- pbus); }
    \coordinate (r2y) at ($(c5.north)+(0,3mm)$);
    \draw (en.south |- pbus) -- (en.south |- r2y);
    \draw (c5.north |- r2y) -- (c8.north |- r2y);
    \foreach \c in {c5,c6,c7,c8}{ \draw (\c.north) -- (\c.north |- r2y); }

    % ---- Recursive composition: Container owns ElementNode children ----
    \draw[compose] (c1.west) to[out=160, in=210, looseness=1.2]
      node[pos=0.12, below left=-0.5mm and 0mm, font=\tiny, align=right]{children 0..*} (en.south west);
  \end{tikzpicture}%
  }
  \caption{Class diagram of the Document Model. Filled diamonds denote composition; the hollow triangle denotes generalisation. \texttt{ContainerNode} composes further \texttt{ElementNode}s, making the document a recursive tree.}
  \label{fig:class}
\end{figure}

\section{Use-Case Diagram}
The system has two actors: the \emph{Author}, a human working in the graphical
editor, and the \emph{MCP Client}, a program or agent driving the same pipeline
headlessly. Figure~\ref{fig:usecase} shows their goals. Note that both
\emph{Author \& export a page} and \emph{Drive builder headlessly} \emph{include}
the shared \emph{Validate document} and \emph{Export ZIP bundle} use cases --- the
same generation, accessibility gate, and packaging serve both the interactive and
the headless path.

\begin{figure}[H]
  \centering
  \resizebox{0.98\textwidth}{!}{%
  \begin{tikzpicture}
    % ---- Use cases: left column ----
    \node[usecase] (uc1) {Author \& export\\a page};
    \node[usecase, below=6mm of uc1] (uc2) {Draw to create\\an element};
    \node[usecase, below=6mm of uc2] (uc3) {Match my\\layout};
    \node[usecase, below=6mm of uc3] (uc5) {Save / open \&\\recover project};
    % ---- Use cases: right column ----
    \node[usecase, right=26mm of uc1] (ucA) {Insert / edit\\elements};
    \node[usecase, below=6mm of ucA] (ucB) {Manage tokens\\\& theme};
    \node[usecase, below=6mm of ucB] (ucV) {Validate\\document};
    \node[usecase, below=6mm of ucV] (ucE) {Export ZIP\\bundle};
    \node[usecase, below=6mm of ucE] (ucH) {Drive builder\\headlessly};
    % ---- System boundary ----
    \node[draw, thick, rounded corners, fit=(uc1)(uc5)(ucA)(ucH),
      inner sep=7mm, label={[font=\bfseries]above:Draw-to-Web}] (sys) {};
    % ---- Actors ----
    \coordinate (author) at ($(uc3.west)+(-40mm,10mm)$);
    \pic at (author) {ucactor};
    \node[font=\footnotesize, below=15mm of author] {Author};
    \coordinate (mcp) at ($(ucH.east)+(40mm,0mm)$);
    \pic at (mcp) {ucactor};
    \node[font=\footnotesize, below=15mm of mcp, align=center] {MCP\\Client};
    % ---- Associations (Author) ----
    \foreach \u in {uc1,uc2,uc3,uc5,ucA,ucB}{ \draw[ucedge] (author) -- (\u.west); }
    % ---- Associations (MCP client) ----
    \draw[ucedge] (mcp) -- (ucH.east);
    % ---- Includes ----
    \draw[ucinclude] (uc1) -- (ucV) node[midway, above, font=\tiny, sloped]{\guillemotleft include\guillemotright};
    \draw[ucinclude] (ucH) to[out=150,in=270] node[pos=0.5, right, font=\tiny]{\guillemotleft include\guillemotright} (ucE.south);
    \draw[ucinclude] (ucH) to[out=120,in=250] (ucV.south);
    \draw[ucinclude] (uc1) to[out=-20,in=160] node[pos=0.5, below, font=\tiny]{\guillemotleft include\guillemotright} (ucE.west);
  \end{tikzpicture}%
  }
  \caption{Use-case diagram. The Author drives the graphical goals; the MCP Client drives the headless goal. Both the interactive and headless export paths \guillemotleft include\guillemotright{} the same validation and packaging.}
  \label{fig:usecase}
\end{figure}

\section{Sequence Diagrams}
Two interactions best characterise the system. The first,
Figure~\ref{fig:seqdraw}, is \emph{draw-to-create}: how a rectangle drawn on the
canvas becomes a real element through an ordinary operation, so that the drawn
pixels themselves never enter the store. The second, Figure~\ref{fig:seqexport}, is
the \emph{export pipeline}: how a document is validated, generated, gated for
accessibility, optimised across the IPC boundary, and packaged. The nine-stage
pipeline is also shown as a flow in Figure~\ref{fig:exportflow}.

\begin{figure}[H]
  \centering
  \resizebox{\textwidth}{!}{%
  \begin{tikzpicture}
    % participants
    \node[seqbox] (p1) at (0,0)    {Author};
    \node[seqbox] (p2) at (37mm,0) {Canvas};
    \node[seqbox] (p3) at (78mm,0) {\texttt{interpret-}\\\texttt{Rectangle}};
    \node[seqbox] (p4) at (119mm,0){\texttt{snapToGrid}};
    \node[seqbox] (p5) at (156mm,0){Store};
    % lifelines
    \foreach \p in {p1,p2,p3,p4,p5}{ \draw[lifeline] (\p.south) -- ($(\p.south)+(0,-102mm)$); }
    % message y-levels
    \foreach \i/\y in {1/-14,2/-24,3/-33,4/-44,5/-54,6/-63,7/-74,8/-83,9/-94}{
      \coordinate (m\i) at (0,\y mm); }
    % messages
    \draw[seqmsg] (p1 |- m1) -- (p2 |- m1) node[midway, above, font=\scriptsize]{drag rectangle};
    \draw[seqmsg] (p2 |- m2) -- (p3 |- m2) node[midway, above, font=\scriptsize]{normalised rect};
    \draw[seqret] (p3 |- m3) -- (p2 |- m3) node[midway, above, font=\scriptsize]{ranked kind + confidence};
    \draw[seqmsg] (p1 |- m4) -- (p2 |- m4) node[midway, above, font=\scriptsize]{confirm kind (picker)};
    \draw[seqmsg] (p2 |- m5) -- (p4 |- m5) node[midway, above, font=\scriptsize]{bounds, grid, breakpoint};
    \draw[seqret] (p4 |- m6) -- (p2 |- m6) node[midway, above, font=\scriptsize]{column start/span + index};
    \draw[seqmsg] (p2 |- m7) -- (p5 |- m7) node[midway, above, font=\scriptsize]{dispatch insertElement(op)};
    \draw[seqret] (p5 |- m8) -- (p2 |- m8) node[midway, above, font=\scriptsize]{updated document};
    \draw[seqret] (p2 |- m9) -- (p1 |- m9) node[midway, above, font=\scriptsize]{element placed};
  \end{tikzpicture}%
  }
  \caption{Sequence diagram: draw-to-create. The canvas interprets the rectangle and snaps it to a grid placement (pure functions), then dispatches a normal \texttt{insertElement} operation; the drawn pixels never enter the store.}
  \label{fig:seqdraw}
\end{figure}

\begin{figure}[H]
  \centering
  \resizebox{\textwidth}{!}{%
  \begin{tikzpicture}
    % participants
    \node[seqbox] (p1) at (0,0)    {Author / UI};
    \node[seqbox] (p2) at (40mm,0) {\texttt{exportProject}};
    \node[seqbox] (p3) at (85mm,0) {Generator\\+ SEO};
    \node[seqbox] (p4) at (125mm,0){\texttt{axeGate}};
    \node[seqbox] (p5) at (162mm,0){Main (IPC)};
    % lifelines
    \foreach \p in {p1,p2,p3,p4,p5}{ \draw[lifeline] (\p.south) -- ($(\p.south)+(0,-110mm)$); }
    % message y-levels
    \foreach \i/\y in {1/-14,2/-24,3/-33,4/-44,5/-53,6/-64,7/-73,8/-84,9/-93,10/-104}{
      \coordinate (m\i) at (0,\y mm); }
    \draw[seqmsg] (p1 |- m1) -- (p2 |- m1) node[midway, above, font=\scriptsize]{export(document, options)};
    \draw[seqmsg] (p2 |- m2) -- (p3 |- m2) node[midway, above, font=\scriptsize]{validate + generate};
    \draw[seqret] (p3 |- m3) -- (p2 |- m3) node[midway, above, font=\scriptsize]{html / css / js};
    \draw[seqmsg] (p2 |- m4) -- (p4 |- m4) node[midway, above, font=\scriptsize]{check(html)};
    \draw[seqret] (p4 |- m5) -- (p2 |- m5) node[midway, above, font=\scriptsize]{violations};
    \draw[seqmsg] (p2 |- m6) -- (p5 |- m6) node[midway, above, font=\scriptsize]{optimize images + minify};
    \draw[seqret] (p5 |- m7) -- (p2 |- m7) node[midway, above, font=\scriptsize]{optimized assets};
    \draw[seqmsg] (p2 |- m8) -- (p5 |- m8) node[midway, above, font=\scriptsize]{bundle ZIP + save(path)};
    \draw[seqret] (p5 |- m9) -- (p2 |- m9) node[midway, above, font=\scriptsize]{ExportResult};
    \draw[seqret] (p2 |- m10) -- (p1 |- m10) node[midway, above, font=\scriptsize]{success / structured error};
    % a11y note (placed in whitespace right of the gate, above the minify message)
    \node[draw, fill=yellow!15, font=\scriptsize, align=center, rounded corners]
      at (147mm,-42mm) {critical / serious\\blocks export};
  \end{tikzpicture}%
  }
  \caption{Sequence diagram: the export pipeline. The accessibility gate runs before any asset crosses the IPC boundary; a critical or serious violation aborts the export with a structured error.}
  \label{fig:seqexport}
\end{figure}

\begin{figure}[H]
  \centering
  \includegraphics[width=0.92\linewidth]{03-export-pipeline.png}
  \caption{The nine-stage export pipeline as a flow: validate $\rightarrow$ generate $\rightarrow$ inject-SEO $\rightarrow$ a11y-gate $\rightarrow$ optimize-images $\rightarrow$ minify $\rightarrow$ sitemap/robots $\rightarrow$ bundle $\rightarrow$ save.}
  \label{fig:exportflow}
\end{figure}
